\honest{Ends Don't Justify Means (Among Humans)}{Eliezer Yudkowsky}{2008}

``I think of myself as running on hostile hardware,'' said Justin Corwin. Yesterday I suggested that humans may have evolved to seek power as idealists and then be corrupted by it, not by plan but by the echo of ancestors who did the same and reproduced: we think we do X for a prosocial reason Y, while other adaptations promote a self-benefiting Z. \nb{That post, ``Why Does Power Corrupt?'', is not in the book.} Today I ask a question well outside classical decision theory: ``What if I'm running on corrupted hardware?''

If so, the feeling that seizing power would help the tribe ``may not be be much evidence'' that it would. The bias does not feel like a bias; its output looks like the way things are. So the human remedy is a rule: ``For the good of the tribe, do not cheat to seize power even for the good of the tribe.'' It says ``even when it would'' help, not ``seems to,'' or people answer that it does not just seem so, it would help if they were in charge. Or do not murder, even for the good of the tribe. \nb{Mill had met this worry in 1861, as an objection to utilitarianism: that a utilitarian ``when under temptation, will see an utility in the breach of a rule.''}

Then comes the philosopher with a train, five people on the track who cannot be warned, and one person I could push in front of it, the only option, and certain to work. An altruistic human's prohibitions ``seem well justified by some historical statistics.'' I give none.

I offer a reply I have ``yet to hear'': running on corrupted hardware, I cannot occupy the epistemic state of knowing those facts with certainty; in a society of AIs without a tendency to be corrupted by power, it would be right for an AI to push the one, and its peers would agree, but I refuse to extend the answer to myself, since that state exists only among other kinds of people. \nb{In R. M. Hare's \textsc{Moral Thinking} (1981), an ``archangel'' with superhuman knowledge and no weaknesses, who needs no moral rules, plays much the same part.} Then I call this reply a dodge, because the universe can force such choices on us, and every legal system decides how many innocent people it will jail to catch the guilty.

My own view: as a human I keep the prohibitions humans made to live in peace, though I do not think them terminally right regardless of consequences, but I would not hold a society of well-calibrated AIs to them, and a Friendly AI that pushed the one person would be no alarm sign (one AI among humans must also consider whether humans learn from its example). I would expect a decent superintelligence to find a better third alternative; but if those are the only two and it judges pushing wiser, even counting the effects on witnesses, I do not object. I do not push people in front of trains or rob banks for altruistic projects; I happen to be human. For a Friendly AI to be corrupted by power would be like it bleeding red blood. The reason is that ``The tendency to be corrupted by power is a specific biological adaptation, supported by specific cognitive circuits''. \nb{In the post's first paragraph humans only ``may have evolved'' this tendency, and the previous day's post had said ``Call it a just-so story if you must.'' No evidence is given for the circuits.}

Minds biased the other way, overestimating the harm of self-benefiting acts, would need the opposite rule, ``the ends do not prohibit the means,'' or they would refuse to breathe for fear of using others' oxygen and all die; for them an occasional selfish overshoot would be as cautiously virtuous as a human passing up a loaf of bread that would really help more than it cost the merchant. So ``the end does not justify the means'' is consequentialism one meta-level up: reflective consequentialism, ``for beings who know that their moment-by-moment decisions are made by untrusted hardware.'' \nb{This is the standard indirect consequentialist view, held in various forms by Mill, Sidgwick and Hare. Its best-known difficulty, that a person who knows the rules are only tools may not hold them firmly, is taken up in the next post, ``Ethical Injunctions.''}
